% Proposed full-depth accessible conceptual-article insertion; no integration.
% Inherit the selected master's sans-serif text and mathematics setup.
\section{A larger model still owes us a learning experiment}

A computer could become better at predicting where a hand will move while the person learns nothing new. That is why the proposed Large Movement Model needs two kinds of test. One asks whether broader training improves a specified prediction. The other asks whether using that prediction helps a trainee, and whether anything learned remains after the teaching commands stop. A better answer to the first question does not settle the second.

The programme's new working hypothesis is that increasing both the amount and the variety of qualified neural, muscle and movement data could help a model cope with new situations. Amount and variety need separate comparisons. Neither has produced a scaling result here, and the proposed body-wide world model has not been trained or validated.

\subsection{Four training sources, five component roles}

The broader proposal draws on at least four training-source groups. Neural observations form one group. Muscle and other myographic records form another; the actual kind of muscle measurement must always be named. Measured motion and displacement form a third. Examples generated by physics-based motion or musculoskeletal simulation form a fourth. A complete collection of each has not been acquired here, and the list leaves room for other useful information.

Suitable neural recordings could come from the brain, spine, peripheral nerves and relevant visceral sources. Their sites, instruments and task roles matter. The present emphasis is motor--sensory. The programme uses Myoception for sensing muscle-driven movement and the feedback that comes with it, including body-position and movement feedback, touch and other task-relevant sensations. This proposed programme use of the name establishes neither a new external scientific standard nor a simultaneous recording of the whole body.

Touch, body-position and movement feedback (proprioception), contact, force and task context each need the right observation. Knowing which tactile stimulus was delivered differs from recording a sensory nerve response or asking what the person felt. A contact sensor's event differs from a calibrated force measurement. Tracking a limb can help estimate its configuration without measuring how that position feels to the subject. EMG records muscle electrical activity; it does not directly record peripheral nerve or sensory activity. Measurements, reports, estimates and simulated values must keep those identities.

The fourth group makes the virtual body a source of training examples as well as a constraint on a forecast. A suitable simulation could generate trajectories, changes between states, calculated muscle--tendon and contact forces, or movement patterns used as a starting point for learning. Record which simulator and body-model version produced them, its parameters, numerical solver, controller and original sources. A controller requesting joint positions and a controller driving muscles do different jobs; the first must never be described as the second.

If real movement is transferred to a different body model, preserve the original recording and the transformation. The internal states or forces calculated from it still belong to the simulation. A generated force is not an independent force measurement. The generator needs calibration against suitable measured support, with uncertainty stated and adjustable settings fitted from training support rather than independent tests.

A proposed test would keep the real training data fixed and add simulation examples in one arm. The permitted information, computing work and original decision rules would be matched, then held-out real observations would check whether the added examples helped. Record how often each arm sees real and simulated examples and how many updates it makes; the same real recordings can still be presented different numbers of times. Count the amount and variety of simulation separately from real acquisitions. Those real tests cannot also be used to calibrate the generator, tune its dynamics, choose rewards or select the model facing them.

The simulation may pass on a useful relationship or its own mistakes in contacts, dynamics and correspondence between bodies. The comparison must be able to expose both. Generated examples cannot independently verify the biological assumptions that produced them. The methods, implementation and use conditions of any proposed generator remain to be checked; no frontier implementation or dataset is selected here.

The four source groups and the programme's five component roles answer different questions. Brain signals, muscle signals, measured movement, a virtual body and a learning model say what the parts do. The four groups say where training examples could come from. The virtual body can supply both examples and physical constraints, and the learning model can use several groups. Additional useful observations remain possible. None of these resources supplies a retained-learning result before the trainee has been tested.

\subsection{Count more examples without confusing them with more experience}

One long recording can supply thousands of short windows. Showing those windows repeatedly gives the model more training presentations, but no additional person, animal or original recording. The proposed study would keep track of those differences: original acquisition groups, subjects, recording time, trials, available target values, distinct windows and repetitions. Copies and overlapping extracts from the same acquisition would stay together when data is divided for learning and testing.

Variety also needs to say what varies. Another person, another recording day, another movement and another sensor each change the problem differently. Body region, organism, apparatus, environment and measurement quality matter too. Counting each collection as one dataset would hide these differences. The study would describe them rather than assume that a larger dataset count means a more informative collection.

An amount comparison would add eligible training support from a declared source distribution. A diversity comparison would keep a specified amount comparable while changing named aspects of its range. Crossing those comparisons could test whether more data helps differently when the sources are broader. The model's design, permitted information, prediction target, preprocessing and selection rules would stay fixed. The study would report repeated exposure, updates to the model and computing work actually used. Equal numbers of updates need not mean equal cost or equal information.

If broader sources also have different measurement quality or easier tasks, that is a combined change. Its outcome cannot be assigned to diversity alone. Exact recordings, budgets, random starting conditions, controls, uncertainty methods and decision rules would have to be fixed before the comparison. This proposed expansion does not change the already adopted continuous WAY-EEG-GAL development experiment or open any original protected test.

\subsection{Let the curve be wrong}

One possible relationship between amount and error is
\[
E(N)=E_{\infty}+A(N/N_0)^{-\alpha}.
\]
$N$ is the chosen training amount and must be greater than zero. $N_0$ is a fixed positive reference amount in the same units, so their ratio is a pure number. $E(N)$ is a nonnegative error on the stated held-out population. In this proposed decreasing family, the floor $E_{\infty}$ and the changing part's size $A$ must each be zero or positive, and the exponent $\alpha$ must be positive. The floor and $A$ have the error's units; the exponent has none. Positive $A$ gives a strictly falling curve as training amount increases. Zero $A$ gives a flat line even when the exponent is positive: the exponent alone does not settle the direction. If the error measure also has an upper limit, predictions must respect it throughout the stated fitting and prediction range. A fitted floor would be a parameter of this curve, rather than an established irreducible biological limit.

No exponent has been fitted here. The measured error is free to stay flat, plateau, rise, reverse or change differently across tasks. A constrained decreasing curve is one candidate explanation; it cannot rule negative transfer out of the results. Those alternatives need to be taken seriously before the results arrive. The study would set its fitting range, competing curve shapes, uncertainty calculation and checks on withheld training-amount levels in advance. A fitted curve must predict data levels left out of fitting. Many overlapping windows from one recording do not become independent repetitions because there are many of them.

Errors also keep their original meaning. An error in a processed muscle signal is different from a camera-pixel error or a joint-angle error. They cannot be mixed into one percentage of biological accuracy without a justified, independently checked scoring rule. Report each measure's units, observation counts, missing values, domains and uncertainty. A bigger computational model is another change again; it cannot be quietly bundled into a claim about more training data.

To test generalization, name what is new. An unseen day, an unseen person, an unseen movement, a new dataset and another species require different comparisons. Each proposed transfer would face learning from target data alone, an appropriate deliberately wrong correspondence and simpler predictors, along with the original tests for shuffled associations, shared cues, recording artefacts and future-information leakage. Poorer performance on particular outputs or domains stays visible. A better average cannot undo an original rule that no evaluated output may be harmed, nor erase exposure during earlier training.

\subsection{Keep the missing parts of a recording missing}

FALCON M1-A offers one candidate connection: motor-cortex activity recorded together with processed muscle electrical signals from electrodes within the muscles of one rhesus macaque. This is intramuscular EMG, which differs from a surface recording over the skin. The recording location and the later signal processing are separate facts. The reviewed release does not establish synchronized motion-capture trajectories. Its muscle targets use processing that includes later information, which limits what an offline result could say about a purely causal real-time prediction. Its small minival subset comes from calibration data and is used to check submission format. Those reused examples do not supply independent validation. \href{https://doi.org/10.48324/dandi.000941/0.241029.1405}{FALCON M1-A release}; \href{https://papers.nips.cc/paper_files/paper/2024/file/8c2e6bb15be1894b8fb4e0f9bcad1739-Paper-Datasets_and_Benchmarks_Track.pdf}{FALCON methods}.

AJILE12 offers another connection: intracranial recordings and upper-body positions estimated from video in twelve clinical participants. These positions are two-dimensional camera pixels. The reviewed source does not establish paired EMG or instrumented three-dimensional motion capture. Its pose estimator is adapted using each participant's labelled video. Whether a particular day later reserved for testing was used to train that estimator remains unverified, as does a numerical bound on alignment error. Labels may therefore be participant-adapted even when the neural prediction model has never trained on that participant. \href{https://doi.org/10.1038/s41597-022-01280-y}{Peterson et al., 2022}; \href{https://doi.org/10.48324/dandi.000055/0.220127.0436}{AJILE12 release}.

The two sources could contribute useful, different training relationships after qualification. They never become the same event. A macaque's brain--muscle recording and a person's brain--pose recording do not supply an observed brain--muscle--movement triple. A model might learn patterns separately, then try an explicit transfer to another source. Actual target measurements would have to test whether that correspondence works.

The release descriptions identify possible source relationships. They do not establish downloaded recordings or a trained model. Exact versions, use conditions, file integrity, original acquisition relationships, timing, units, earlier exposure and the independence of the chosen partitions still need checking before training. Mapping motion between bodies also needs declared geometry, axes, proportions, load, contact, timing and uncertainty. Bodies may share useful mechanical relationships without sharing identical neural codes.

\subsection{Make the inside explain the outside}

When fingers close around an object, the visible movement depends on what happens within the body model and at the contact. The proposed world model would connect specified joint and muscle mechanics to supported movement, objects and environmental conditions. Its interior initially means those named mechanical states. It does not mean that every organ, nerve or sensation has been reconstructed.

NeuroTranslate would map named observations to specified targets. MoCap ReComp would supply a separately qualified movement composite that retains where each observation came from. An ordinary set of estimated video poses does not become that proposed instrument by receiving its name. MyoSim Replay would calculate the declared body's mechanical response. Their proposed combination is the Large Movement Model. Bringing neural representations from different recording types into one model would not establish matching individual cells across organisms or assemble a measured whole nervous system.

A proposed internal state could describe joint arrangement, generalized velocity, muscle-model activation and the environment. Rotational positions and velocities may use radians and radians per second; sliding coordinates require length units. A transition model would use permitted earlier observations, context and declared controls to forecast possible later states and their uncertainty. Every instrument would have its own observation relationship. EMG supplies muscle electrical activity, not direct peripheral nerve activity, and a processed EMG amplitude does not automatically measure activation or force. A tracked position is different from an internal force calculated by the simulator. Missing observations need an availability mask, and numerical filling values must never acquire the status of measured zeros. Predicting recorded neural features from body observations can leave several hidden explanations possible.

The physics needs more than a believable picture. Predicted movement must be checked against independent measured motion in matching time, coordinates and units. A claim about internal force needs suitable force evidence as well: a similar path can be produced by different combinations of muscle forces. Joint limits, loads, contacts, gravity, friction, energy accounting and sensitivity to numerical settings need their own checks. OpenSim offers established methods for specifying musculoskeletal models; it does not validate this proposed assembly. \href{https://journals.plos.org/ploscompbiol/article?id=10.1371/journal.pcbi.1006223}{Seth et al., 2018}.

% R4 full-depth accessible conceptual-article insertion; no integration.
% Inherit the selected master's sans-serif text, code and mathematics fonts.
\subsection{One muscle pulse, two simulated trajectories}

A computer hand can respond to a muscle command before any person has learned a skill. This test asks about that response. We compared two trajectories of \texttt{myohand\_r}, an installed mechanism representing a human right hand attached to a fixed passive torso. Both started in the same recorded state; one received a short muscle excitation. The comparison tests whether the specified command changes the state the software computes. It contains no brain decoding, measured biological response or evidence of retained learning. The software was \texttt{myo\_sim} 0.2.3 at commit \texttt{93b0ca8f4ec90c9899ee7f05fee561e9911da91b}, MuJoCo 3.14.0 and NumPy 2.5.3.\cite{myohandMyoSimPinned2026,myohandMuJoCo3140,mujoco2012,myohandNumericalRecord20261105}

\subsubsection{Let the control move too}

The compiled hand has 23 scalar hinge coordinates, meaning 23 angles describing its movable joints, measured in radians. It has 39 muscle--tendon actuators that calculate muscle-driven pulls, with 39 internal activation states and zero sensors. These are parts of a computer mechanism. Its activation values are not measured muscle electrical signals, and its forces are not readings from a force instrument.

The two branches began with matching complete recorded initial states, including angles, velocities and muscle-model activation. One received no additional excitation. The other received dimensionless excitation 0.2 at the actuator named \texttt{FDS2\_r}; all other control values were zero. Dimensionless means that 0.2 uses the model's input scale, rather than volts, newtons or an EMG amplitude. It drives a muscle model instead of requesting a joint position directly.

Both branches used the native Euler setting, the simulator's specified rule for advancing its state. Each advanced 100 times by $0.002\,\mathrm{s}$, reaching nominal $0.200\,\mathrm{s}$. Starting the count at step zero, the pulse ran at steps 10 through 59. Its pre-step interval was $0.020\leq t<0.120\,\mathrm{s}$. Gravity and passive mechanics could move either hand. The unpulsed branch therefore supplies a moving comparison, rather than a requirement that the hand freeze when the added command is absent.

The state-dependent quantities were recalculated with a forward evaluation at the start and end of each branch. The saved record contains 200 explicit integration steps and four explicit forward evaluations. The executed construction path also implies two compilations and one internal forward evaluation while the mechanism was being built. These are counts of identified calls, not a tally of everything inside the native solver.

The construction settings need care. The optional \texttt{inertia\_floor} argument stayed at \texttt{None}, yet the native builder changed each diagonal rotational-inertia component of four fixed scaffold bodies from $1$ to $10^{-4}\,\mathrm{kg\,m^2}$. Rotational inertia describes resistance to changes in turning motion. That change belongs to the tested mechanism. No extra physical parameters were retuned and no alternative \texttt{implicitfast} fallback was used. This check does not calibrate the hand's position, orientation, size, contacts or fidelity to biology.

\subsubsection{Compare the two endpoints}

Use $P$ for the pulsed hand, $C$ for the control and $T=0.200\,\mathrm{s}$ for the nominal endpoint. The three reported differences are
\[
\Delta q_{\infty}=\max_{1\leq j\leq23}|q_{P,j}(T)-q_{C,j}(T)|,
\qquad
\Delta a=a_{P,\mathrm{FDS2}}(T)-a_{C,\mathrm{FDS2}}(T),
\]
\[
\Delta F=F_{P,\mathrm{FDS2}}(T)-F_{C,\mathrm{FDS2}}(T).
\]
The first calculation finds the difference between corresponding final joint angles, discards its sign and selects the largest of the 23 values. It describes the final instant, rather than the largest difference anywhere along the paths. The second subtracts the control activation from the pulsed activation for the selected muscle model. The third makes the same signed subtraction for its tendon-actuator force. Activation has no physical unit; force uses newtons and the model's sign convention. A forward evaluation recalculated each final force at its corresponding final state, avoiding force values left over from the last step.

% Shared by all four proposed R3 inserts; no added packages or font defaults.
\begin{table}[htbp]
\centering
\caption{Final contrasts for one simulated right-hand mechanism. All values compare the pulse branch with its matched control at nominal $T=0.200\,\mathrm{s}$. The angle value is the largest absolute separation across 23 hinges at that endpoint; activation and force retain the pulse-minus-control sign. Activation and force concern the selected \texttt{FDS2\_r} actuator. Force uses the synchronized final forward evaluations and the model's sign convention. These are simulated contrasts, not errors against biological measurements. The displayed digits preserve the recorded numerical values, without establishing biological accuracy.\cite{myohandNumericalRecord20261105}}
\label{tab:myohand-numerical-response}
\begin{tabular}{p{0.53\linewidth}p{0.26\linewidth}p{0.10\linewidth}}
\hline
Quantity & Contrast & Unit \\
\hline
Largest absolute final joint-coordinate separation, $\Delta q_{\infty}$ & $0.5027347840656727$ & rad \\
Selected final activation difference, $\Delta a$ & $0.05205729904534248$ & none \\
Selected final tendon-actuator force difference, $\Delta F$ & $-6.407722136774943$ & N \\
\hline
\end{tabular}
\end{table}


A separate standard-library audit checked the saved evidence without rerunning the model. It completed 132 checks without a failed check. The required numerical channels were finite, and recorded native warning counts stayed at zero. Four saved endpoint mass-matrix spectra contained positive values. A mass matrix describes the model's resistance to acceleration; its spectrum is a numerical diagnostic of that matrix. The raw matrices and the inputs to residual calculations were not independently exported. The audit therefore checked the stored diagnostics without independently solving those matrices again.

There was no observed contact. At the final control endpoint, the recorded soft-limit excess was $0.0004121829780273552\,\mathrm{rad}$. A soft limit resists crossing a boundary without proving an absolute hard stop. Neither energy conservation nor numerical convergence was established. The test imposed no physiological threshold for calling the response large enough.

The result is a difference between two simulated trajectories. It is not an error measured against a person's movement, muscle signals or forces. The 200 time steps depend on earlier states; they do not supply 200 subjects or 200 independent experiments. No statistical $p$-value follows. The decimal places preserve what the numerical record contains, without establishing biological accuracy or repeatability in other cases.

\subsubsection{Follow the evidence through to the trainee}

The diagnostic establishes the recorded input-to-state response, not a performance benchmark. Its complete computing-resource record and unresolved memory/timing attribution are retained in Appendix~\ref{app:reproducibility-availability}.

With zero sensors, the hand records none of Myoception's proposed sensory side of muscle-driven action, such as touch or body-position feedback. Calculated activation, force and motion are different from sensory nerve recordings or measured biological feedback. This test trains no controller or Large Movement Model and tests no generalization. The wider proposal still draws on four groups: neural observations, muscle measurements, measured movement and physics-generated examples. Before this generator could supply training examples, its mechanics, controls, calibration, uncertainty, source rights and data lineage would need qualification. Independent held-out real measurements would then test its contribution. More data and more varied data remain separate, untested hypotheses.

Even a later controller that passes its mechanical tests would leave Brain2Body's teaching question open. A trainee must be tested after trainer commands end, with any necessary compensatory movement bridge kept constant so that movement can still be expressed. Match the information, practice and assistance, retain the original sham, delayed and unrelated-guidance controls, and compare independent outcomes. No trainee, teaching, withdrawal or physical intervention appears in this diagnostic. It chooses no training generator and supplies no biological ground truth. LINK's eighteen fits and 720 saved epochs still fail both admission gates and all five benefit claims; protected tests and quarantined sources remain closed.

\subsection{Follow the gain until the teaching ends}

A better movement prediction might make assistance more useful. To test teaching, the experiment must also identify what the trainer supplied: expert information, feedback, commands, practice and the amount of help. Match those resources against an appropriate simpler controller and the original controls using sham, delayed or unrelated guidance. Otherwise, extra practice or being given the correct action could explain the apparent benefit.

The retention test stops trainer commands while keeping any necessary compensatory movement bridge constant. That bridge lets a trainee express movement despite an interrupted pathway; removing it together with teaching could make movement impossible regardless of what was learned. Immediate performance, later retention and unfamiliar variations are separate outcomes. A gain that requires continued guidance can support an assistance claim. A gain that persists under the controlled comparison needs its own uncertainty and tested domain. It still would not establish tissue repair or the ability to transfer any skill to any person.

Offline replay, simulation, measured biological feedback and physical intervention keep their own evidence and authority. The current proposal reports no scaling result, trained body-wide world model or biological teaching. LINK's eighteen fits and 720 saved epochs retain their failed model-admission gates and all five failed benefit claims. The next comparison has to show which prediction improves, what supplied the improvement and whether the trainee retains anything from it.
